\documentclass[UTF8,zihao=-4,fontset=windows]{ctexart}
\usepackage[a4paper,margin=24mm]{geometry}
\usepackage{booktabs,tabularx,graphicx,xcolor,fancyhdr}
\usepackage[hidelinks]{hyperref}
\definecolor{ink}{HTML}{173E52}
\ctexset{section={format=\large\bfseries\color{ink}}}
\setlength{\parskip}{6pt}\setlength{\emergencystretch}{3em}
\pagestyle{fancy}\fancyhf{}\fancyhead[L]{\small 阶段报告 6 · 中文导读}\fancyfoot[C]{\thepage}\setlength{\headheight}{15pt}
\title{\bfseries 模型能记住训练图，\\为什么换一张图还是找不全？\\[8pt]\large 阶段报告 6 导读}
\author{Xinhao Jin · bdd26}\date{2026 年 10 月 26 日}
\begin{document}\maketitle

\section*{这两周主要做了什么}

上一阶段，我们检查过“找到了核，却分错类别”的问题。简单地换分类器没有带来稳定改善，也看到了一些标签边界可能不清的证据。从 10 月 13 日起，研究进一步往前移：先不急着让模型分对类别，而是检查它能不能真正找到一组完整、可信的细胞核。

这段时间试了三类办法：改变裁切后小核的训练监督；用查询模型一次生成一组独立实例；先找中心点，再考虑如何画出轮廓。完整训练和小规模诊断都已积累了不少结果，但目前还没有一个新配方通过原先定好的条件，可以进入最终大规模确认。

这不是“没有任何数字上涨”。恰恰相反，多次出现了某项指标变好、另一项变差的情况。本阶段的重要工作，是看清这些上涨究竟意味着什么，避免把更高的训练拟合、更少的预测，或者真值帮助下的理想结果误当成方法进步。

\begin{quote}\small\textbf{本轮核实状态：}P16 seed 19 两组已完成并核验；seed 1 仍运行。P17 真实 V100 验证 55297045 已通过；正式数组 \textbf{55297061} 于 2026 年 10 月 26 日 \textbf{09:08:19（美东时间）}单次提交，四项初始均排队。提交不代表训练完成，也不代表方法通过。\end{quote}

\section*{第一件事：发现弱点，不代表已经知道怎么修}

我们把原本完整的核裁成只露出一部分，发现模型确实更容易漏掉它们。这个现象在修正推理模式、考虑面积差异后仍然成立。于是进一步改变训练监督，希望模型学会处理这类局部可见的小核。

但四组、两个随机种子的完整训练没有支持这套修复办法。普通裁切增强让 Dead 类的某些指标提高，却降低了整体分割质量；加上选择性监督以后，表现也没有进一步变好。在固定裁切样本中，两种子的效果还一正一负。

所以我们保留“裁切是一个脆弱点”，停止“这一套监督已经能修复它”的设想。前一个观察是研究线索，后一个需要真正的干预实验，不能跳过。

\section*{第二件事：画出一些像样的轮廓，还不等于找到了一组核}

查询模型的想法是：让很多查询各自负责一个对象，再把它们组成最终结果。实际训练却出现了严重退化。把学习率降得慢一些，也没有救回来。以同轮基线约 .473 的 mPQ 为参照，查询模型只有约 .088；这个差距不是小幅随机波动。

进一步查看同一批训练图，问题更清楚了。331 个目标中，一个种子的原始查询掩码能与 212 个目标有足够重叠；模型自己筛选完查询，只剩 104 个；处理轮廓重叠后剩 84 个。另一个种子是 227、119、97。

这说明模型并非完全没有生成可用轮廓，而是有一大部分没有被保留成最终实例。同时，原始掩码也远没覆盖全部目标。因此，单纯降低一个置信度阈值不够，继续扩大同一套训练也缺少依据。要分别检查：它有没有画出对象、有没有把它当成对象、对象之间冲突时有没有丢掉它。

\section*{第三件事：有标准答案帮助时能做，不等于推理时会做}

有些小实验直接给模型真值中心或对象对应关系，再让它画轮廓。这时，训练图上可以拟合相当一部分 Dead 核。乍看像是好消息，但标准答案已经替模型完成了最难的一部分——告诉它该看哪里。

再看两套真正预测点的配方：E17 没有产生候选；E18 产生很多候选，但最终实例推理的审计精度只有约 3.5\%--3.7\%。这是实例推理结果，不是后续实验的中心点匹配率。另一项真值点实验中，本轮读出留出图像也明显弱于拟合图像；不过冻结的 E7 已经见过这些图，不能把它说成整模型对新图像的独立检验。因此，我们没有直接把这些模块拼起来训练完整 PM-Net，而是先检查真正预测的候选是否既全又准。

\section*{为什么反复强调“训练图”和“留出图”}

在后续 P3 起的实验中，可以把拟合图理解为平时练习过的题，审计图是没有参与该轮参数拟合的题。一个模型把练习题记住，并不代表它学会了处理同类新题。这个定义不能倒用到前面所有小诊断上，每项实验的留出范围要分别看。

最早的冻结特征点头在补足训练以后，训练内召回已经约 94\%--98\%，但换到开发图像，内部 Dead 点召回仍只有约 40\%。后来允许模型适配特征，并加入增强，训练图表现进一步改善；在留出图像上，却总是在“少漏掉一些”和“别报错太多”之间付出代价。

为了把这种落差看清，后续在训练折内部固定留下 522 张图，另外 2,134 张用于拟合。这个审计集里只有 11 张内部 Dead 阳性图、145 个对象，而且大部分来自 Lung 和 Colon。后面所有点检测配方都用同一批图来作比较，方便控制变量，也意味着不能把一轮轮结果当成新的独立验证。

\section*{连续试验排除了哪些简单办法}

我们分别检查了更细的输出网格、更宽的目标邻域、颜色扰动、稀有类加权、少训练一些骨干层、dropout、稀有类重采样、更小的骨干学习率，以及平移增强。它们都没有同时满足原先的双种子通过条件。

其中两个例子尤其有解释力。把四分之一的训练位置换成稀有类阳性图以后，拟合 Dead 召回几乎 100\%，留出图像却降到约 35\%--36\%。这套方案不能靠“让模型多看几遍稀有类”解决问题。

平移增强是另一种失败。模型在训练图上本身就明显学不好：总体召回从约 98\% 降到 47\%。它报得少了，所以精度从约 69\% 升到 86\%，但正确找回的核也少了很多。只看精度就会得出相反的结论。白色填边、部分核被移出画面和更强的任务要求同时改变，不能简单归咎于某一种位置编码。

\section*{小核确实难，但还不能说原因只有尺寸}

高分辨率适配器也做了完整训练，却没有稳定提高内部 Dead 的匹配率。对四个模型共同失败的对象作检查，发现小核确实占多数；但较大的对象同样会持续失败。

更细查训练数据后，尺寸和对比度各个粗分组里都能找到训练样本。也就是说，不能简单解释为“训练里没有这种小核”。但有相似尺寸的样本，不代表同样的组织、外观和患者差异都覆盖了。现在还不能把这些相关性说成唯一原因，更不能在缺少专家判断时宣布标签一定有错。

\section*{最新结果，以及为什么下一轮换了一个问题}

P16 把只在精确中心给正监督的损失，改成直接拟合完整 Gaussian 热图。已完成的 seed 19 中，Dead 点召回从 48.28\% 提高到 51.72\%，多匹配了 5 个对象；但区间仍跨零，精度约 70.11\%，都不足以通过原条件。另一个种子尚未收齐，目前不能给出完整实验结论。

下一轮 P17 比较两套表示：一套是原 UNI 检测器，另一套是从头训练的卷积金字塔。问题是：在同一批数据、损失、解码和预算下，换一种表示配方是否值得继续？这轮没有同时加掩码头、查询或新的稀有类技巧，也不把它包装成一个已经有效的新方法。卷积模型与对照的预训练和学习率不同，结果只能评价整套配方，不能单独证明“局部特征就是原因”。

\par\textbf{下一轮已经提交。}10 月 26 日 09:08:19（美东时间），四项正式对照以作业数组 \textbf{55297061} 单次提交：两种模型各做两个随机种子，每项 130 轮，同时最多运行两项。此前的真实 V100 验证已完成，291 项检查通过。正式作业提交时仍在排队，尚没有训练结果；提交成功不能替代效果检验。


\section*{接下来怎样判断，而不是继续盲目试}

先把 P16 剩余运行收齐，再核验 P17。每轮都要检查实际训练轮数、来源、配置和保存文件，再按两个种子分别套用原标准。对于点检测，必须既找得全又报得准；只有这样，接掩码模型才有实际意义。

本阶段没有新的三划分测试成绩。图像层面的抽样区间不等于患者层面的保证，多轮使用同一开发数据也会使后续判断更容易乐观。已有负结果应当约束下一步，而不是被换一种说法重新打开。完整数字、指标定义、停止理由和来源见正式报告及核验文件。
\end{document}
